\section{Introduction}
Stars maintain hydrostatic equilibrium for the majority of their lifespan as pressure from fusion balances gravity, but as fuel runs low and fusion processes slow, stellar cores shrink from the inbalance. For stars under eight solar masses ($<8M_\odot)$, the core shrinks until electron degeneracy pressure dominates, forming a white dwarf after the red giant phase, where Chandrasekhar famously found an upper mass limit of the resulting object ($\sim1.4M_\solar$)\cite{Chandrasekhar1935}. The heaviest stars ($\geq 25M_\odot$) leave behind a black hole after going supernova from core collapse as the schwarzchild radius is crossed \cite{Heger2003}. In the middle, we get some of the most extreme places in the universe, the neutron star. Neutron stars offer themselves to be a great observational tool as their special degenerate interiors are impossible to replicate on earth. These incredibly dense stellar objects were first theorized by Baade and Zwicky. They proposed that tightly packed neutrons in a stellar object would be responsible for the immense release of energy from observed supernova (a contender for the origin of cosmic rays). \cite{Baade} The statistical mechanics was then heavily pioneered by Oppenheimer and Volkoff as they explored the structure of neutron stars by treating the interior as a cold and relativistic neutron gas under hydrostatic equilibrium in 1939. \cite{Oppy} Of course, we didn't get any neutron star observations until X-Ray astronomy kicked off in the 1960s and additional improvements to radio astronomy were made, where Giacconi indirectly observed extra-solar X-Rays and speculated sychrotron or stellar origins. \cite{Giacconi} A few years later in 1967, an exciting discovery was made by Jocelyn Bell Burnell, who identified a rapid periodic signal while poring over pages of measurements. Potential intra-solar sources were crossed out and there was wonder of extra-terrestial intelligence, but similar pulsating sources identified across the sky solidifed that it was some natural phenomena. However, Bell Burnell's supervisor, Hewish, would be first author in the following publication, even after having pushed back on the statistical significance of the signal. He would then go on to recieve the Nobel Prize in Physics in 1974 after Gold and Pacini showed that these pulsating sources \textit{were} rapidly rotating neutron stars, pulsars.\cite{apsnews_pulsars_2006,Hewish1968} Bell Burnell's contribution would only be recognized years later when she was awarded the Special Breakthrough in Fundamental Physics award in 2018.\cite{fienberg_bellburnell_2018}. Today, thousands of pulsars have been observed as recorded by the Australian Telescope National Facility Pulsar Catalogue.\cite{manchester_atnf_2005} The highly regular period of the pulsar pulse inspired Carl Sagan and Frank Drake to use this property to pinpoint earth on a pulsar map onboard the voyager probe. \cite{nasa_pia14113_voyager}

As the name suggests, neutron stars are primarily composed of neutrons. They typically have radii of $R\sim\qty{10}{km}$ and masses of about $M\sim 1-3M_\solar$\cite{shapiro_teukolsky_1983}. The implications for the density are immense when considering that matter which is heavier the sun is being compressed into a space smaller than the width of the city of Calgary. This corresponds to an average density of about $\rho \sim \qty{10}{^{15}g/cm^{3}}$, exceeding the nucleon density of BLAH, which from a classical perspective looks as if the neutrons are all touch eaching other.[CITE] From here we must ask ourselves the central question of this paper. \textit{What stops a neutron star from collapsing in on itself?}. The answer will come from neutron degeneracy pressure associated with the pauli exclusion principle as it applies to a fermi gas. For the rest of this paper, we will explore the quantum fermi gas model, the composition and formulation of neutron stars, use non-relativistic and relativistic limits to estimate mass and radii, and then explore the mass-radius relationship before discussing limitations. 